\section{Пространства интегрируемых функций}
\label{sec:integrable-function-spaces}

\subsection[Пространства Lp]{Пространства $L^p$}

Пусть $E\subset\mathbb R^n$ --- измеримое множество. Для
$1\leq p<\infty$ рассматривают множество измеримых функций, для которых
\[
    \int_E |f(x)|^p\,dx<\infty.
\]
Естественная величина
\[
    \|f\|_p=
    \left(\int_E|f(x)|^p\,dx\right)^{1/p}
\]
не различает функции, отличающиеся лишь на множестве меры нуль. Поэтому
элементами пространства $L^p(E)$ являются не отдельные функции, а классы
эквивалентности по отношению
\[
    f\sim g\quad\Longleftrightarrow\quad f=g\quad\text{почти всюду на }E.
\]
На этих классах $\|\cdot\|_p$ является нормой.

Для $p=\infty$ вводится пространство существенно ограниченных функций
\[
    L^\infty(E)=\{f:\ \operatorname*{ess\,sup}_{x\in E}|f(x)|<\infty\},
\]
с нормой
\[
    \|f\|_\infty=\operatorname*{ess\,sup}_{x\in E}|f(x)|.
\]

\subsection{Неравенства Гёльдера и Минковского}

Пусть $1<p<\infty$, а $q$ --- сопряжённый показатель:
\[
    \frac1p+\frac1q=1.
\]
Для $f\in L^p(E)$ и $g\in L^q(E)$ выполняется \emph{неравенство Гёльдера}
\[
    \int_E|fg|\,dx
    \leq \|f\|_p\|g\|_q.
\]
Отсюда, в частности, следует, что произведение $fg$ суммируемо.

Для $f,g\in L^p(E)$ выполняется \emph{неравенство Минковского}
\[
    \|f+g\|_p\leq\|f\|_p+\|g\|_p.
\]
Именно оно даёт неравенство треугольника для нормы $L^p$. Вместе с
однородностью и положительной определённостью это показывает, что $L^p(E)$ ---
нормированное линейное пространство.

В случае конечной меры из неравенства Гёльдера также удобно получать оценки
интегралов через норму функции. Например,
\[
    \|f\|_1\leq \mu(E)^{1-1/p}\|f\|_p,
    \qquad 1<p<\infty.
\]

\subsection[Полнота пространств Lp]{Полнота пространств $L^p$}

Одно из фундаментальных свойств пространств Лебега состоит в их полноте:
\[
    L^p(E)\ \text{--- банахово пространство},
    \qquad 1\leq p\leq\infty.
\]
То есть если последовательность $\{f_n\}$ фундаментальна по норме $L^p$,
\[
    \|f_n-f_m\|_p\longrightarrow0
    \qquad (n,m\to\infty),
\]
то существует $f\in L^p(E)$ такая, что
\[
    \|f_n-f\|_p\to0.
\]
Полнота принципиальна в анализе и численных методах: предельный объект остаётся
в том же функциональном пространстве.

\subsection[Пространство L2]{Пространство $L^2$}

Особое место занимает $L^2(E)$. В нём норма порождается скалярным
произведением
\[
    (f,g)_{L^2}=\int_E f(x)\overline{g(x)}\,dx,
    \qquad
    \|f\|_2=\sqrt{(f,f)_{L^2}}.
\]
Поэтому $L^2(E)$ является гильбертовым пространством. В нём имеют смысл
ортогональность,
\[
    (f,g)=0,
\]
ортогональные проекции и разложения по ортонормированным системам. Именно поэтому
$L^2$ особенно часто возникает в вариационных и спектральных задачах.

\subsection{Сопряжённые пространства}

Для $1<p<\infty$ каждый элемент $g\in L^q(E)$, где
$1/p+1/q=1$, задаёт линейный непрерывный функционал на $L^p(E)$:
\[
    \ell_g(f)=\int_E f(x)\overline{g(x)}\,dx.
\]
По неравенству Гёльдера
\[
    |\ell_g(f)|\leq\|f\|_p\|g\|_q.
\]
Теорема Рисса о представлении утверждает и обратное: любой линейный
непрерывный функционал на $L^p(E)$ имеет такой вид с единственным
$g\in L^q(E)$, причём
\[
    \|\ell_g\|_{(L^p)^*}=\|g\|_q.
\]
Следовательно,
\[
    (L^p(E))^*\cong L^q(E),
    \qquad 1<p<\infty.
\]
В частности, при $p=q=2$ пространство $L^2$ изометрически отождествляется со
своим сопряжённым, что согласуется с общей теоремой Рисса для гильбертовых
пространств.

\subsection{Плотные классы функций}

Для $1\leq p<\infty$ простые функции плотны в $L^p(E)$. При стандартных
условиях на область также можно приближать элементы $L^p$ непрерывными и
гладкими функциями. Это свойство позволяет сводить многие доказательства и
численные построения сначала к более удобным функциям, а затем переходить к
пределу по норме $L^p$.

\subsection{Что важно сказать на экзамене}

Пространство $L^p$ состоит из классов измеримых функций, равных почти всюду,
с нормой
\[
    \|f\|_p=\left(\int|f|^p\right)^{1/p},
\]
а для $L^\infty$ используется существенный супремум. Основные неравенства ---
Гёльдера и Минковского. Все $L^p$, $1\leq p\leq\infty$, полны; $L^2$ является
гильбертовым пространством. Для $1<p<\infty$ сопряжённым к $L^p$ является
$L^q$, $1/p+1/q=1$, и функционалы имеют вид
$\ell(f)=\int f\overline g$ (в вещественном случае черта комплексного сопряжения отсутствует).

\newpage
